\chapter{L'apparato circolatorio}

% ---------------------------------------------------------------------------
\section{Generalità}

\textbf{Cuore}: pompa. \textbf{Vasi sanguiferi}: vie di conduzione del sangue.

L'apparato circolatorio è l'insieme di organi deputati al trasporto dei liquidi: è un sistema
chiuso di vasi, in cui il sangue circola sotto la spinta del cuore.

\rubrica{Caratteristiche dei vasi sanguigni}
\begin{itemize}
  \item \textbf{resistenti}, per sopportare le variazioni della pressione ematica, che varia con
    la distanza dal cuore;
  \item \textbf{flessibili}, per adattarsi ai movimenti degli organi vicini;
  \item \textbf{robusti}, per evitare di essere occlusi da pressioni esterne;
  \item \textbf{adattati nella struttura} alla specifica funzione.
\end{itemize}

I vasi sanguigni si suddividono in: arterie, arteriole, capillari, venule, vene.

\rubrica{Funzioni dei vasi}
\begin{itemize}
  \item trasporto del sangue;
  \item scambio gassoso anidride carbonica/ossigeno nei capillari polmonari, e scambio gassoso
    ossigeno/anidride carbonica e di nutrienti/sostanze di rifiuto nei capillari sistemici;
  \item regolazione della pressione sanguigna;
  \item indirizzamento del flusso sanguigno sistemico verso i distretti che in quel dato istante
    necessitano di un maggiore apporto di sangue: es.\ l'apparato digerente durante la digestione,
    o i muscoli durante l'attività fisica.
\end{itemize}

% ---------------------------------------------------------------------------
\section{Struttura istologica dei vasi sanguigni}

\begin{itemize}
  \item \textbf{tonaca intima}: endotelio monostratificato piatto e sottostante strato di tessuto
    elastico (endotelio + strato subendoteliale fibroso, o membrana basale, + membrana elastica
    interna; quest'ultima più spessa nelle arterie);
  \item \textbf{tonaca media}: strati concentrici di tessuto muscolare liscio con fibre collagene
    (nelle arterie è presente anche una banda esterna di tessuto elastico, la membrana elastica
    esterna);
  \item \textbf{tonaca esterna} (o avventizia): tubo fibroso robusto contenente vasi (\textit{vasa
    vasorum}) e nervi, che ancora i vasi al connettivo circostante.
\end{itemize}

% ---------------------------------------------------------------------------
\section{L'albero arterioso}

L'albero arterioso si struttura in:
\begin{enumerate}
  \item \textbf{arterie}, in cui è fondamentale la componente elastica;
  \item \textbf{arteriole}, in cui è fondamentale la componente muscolare;
  \item \textbf{capillari arteriosi};
\end{enumerate}
a questi fanno seguito: capillari venosi, venule post-capillari, vene.

\rubrica{Arterie elastiche}
Sono le più vicine al cuore, in cui prevale la componente elastica; ricevono una forte spinta
sistolica, ed hanno un calibro elevato ($1$--$3\,cm$). Esempio: l'\textbf{aorta}. Contengono una
notevole quantità di elastina: durante la sistole ventricolare si espandono incamerando energia
potenziale, che restituiscono durante la diastole ventricolare, contribuendo a:
\begin{itemize}
  \item spingere il sangue in avanti;
  \item smorzare le oscillazioni pressorie;
  \item trasformare il flusso da intermittente a continuo.
\end{itemize}

\rubrica{Arterie muscolari}
Prevale la componente muscolare; diametro $0,5\,mm$--$1\,cm$. Esempio: l'\textit{arteria
femorale} (spessa lamina elastica).

\rubrica{Arteriole}
Scompare l'avventizia, la tonaca media è muscolare; diametro $< 0,5\,mm$.

% ---------------------------------------------------------------------------
\section{Capillari}

\textbf{Capillari}: cellule endoteliali + lamina basale; diametro $5$--$10\,\mu m$; lunghezza
media $1\,mm$. Sono i vasi più piccoli, deputati allo scambio coi tessuti: spesso una sola cellula
endoteliale forma tutta la sezione trasversa del capillare.

Sono molto numerosi in tutti i tessuti, con alcune eccezioni: tendini e legamenti; cartilagine;
epiteli; cornea.

Essendo molto numerosi ed avendo una parete molto sottile, rappresentano le strutture ideali per
fornire alle cellule sostanze indispensabili quali gas, nutrienti, ormoni.

\rubrica{Tipi di capillari}
\begin{itemize}
  \item \textbf{continui}: i più comuni nell'organismo, abbondanti nella pelle e nei muscoli. Le
    cellule endoteliali formano un rivestimento continuo; giunzioni serrate mantengono le cellule
    aderenti, lasciando solo piccoli spazi dai quali escono piccole quantità di liquidi;
  \item \textbf{fenestrati}: simili ai continui, eccetto per la presenza di pori che attraversano
    a tutto spessore la cellula. Questi pori sono coperti da sottili diaframmi, e poi dalla lamina
    basale delle cellule endoteliali; le fenestrature consentono un notevole aumento di
    permeabilità. Si trovano preferenzialmente dove devono avvenire in maniera rilevante:
    \begin{itemize}
      \item formazione di ultra-filtrato (rene, plessi corioidei);
      \item riassorbimento di liquidi e sostanze in esso disciolte (digerente, rene);
      \item contatti con l'apparato vascolare: organi endocrini;
    \end{itemize}
  \item \textbf{sinusoidi}: simili ai fenestrati, ma con forma irregolare e tortuosa; pori
    ingranditi, spazi anche ampi tra le cellule endoteliali; lamina basale più sottile; lento
    movimento del sangue. Queste modificazioni consentono sia ai liquidi che a macromolecole e
    cellule del sangue di entrare nello spazio interstiziale o nel lume del capillare. Presenti in
    fegato, milza, midollo osseo ed organi endocrini: con interruzioni intercellulari endoteliali
    e membrana basale incompleta, permettono anche il passaggio di cellule.
\end{itemize}

\rubrica{Edema}
Ristagno di liquido nei tessuti. Cause:
\begin{itemize}
  \item aumento della pressione venosa (insufficienza cardiaca o ristagno venoso);
  \item ristagno del flusso linfatico;
  \item capillari divenuti permeabili alle proteine plasmatiche (infiammazione, risposte
    immunitarie);
  \item diminuzione delle proteine plasmatiche (malnutrizione, malattie epatiche, etc.: ascite).
\end{itemize}

% ---------------------------------------------------------------------------
\section{Vene}

Struttura analoga alle arterie ma con minore componente muscolare ed elastica, e più connettivo.
Anche le vene sono innervate dall'ortosimpatico, che facendole contrarre costringe il sangue a
fuoriuscire da alcuni organi (fegato, milza, polmoni, ad elevata funzione capacitativa) per
incrementare la pressione del sangue.

Dai capillari il sangue passa nelle: venule post-capillari ($10$--$25\,\mu m$)
$\rightarrow$ venule collettrici ($20$--$50\,\mu m$) $\rightarrow$ vene muscolari
($50$--$100\,\mu m$) $\rightarrow$ vene (fino a $4\,cm$).

\rubrica{Confronto vene/arterie}
\begin{itemize}
  \item \textbf{vene}: parete più sottile; lume collassabile; struttura scarsamente elastica;
    possono essere dotate di valvole; bassa pressione;
  \item \textbf{arterie}: parete più spessa; forma cilindrica sempre conservata; struttura molto
    elastica; prive di valvole; alta pressione.
\end{itemize}

Nella maggior parte dei casi le vene hanno un decorso speculare a quello delle arterie, di cui
assumono il nome e che accompagnano nel decorso.

Le vene degli arti e della parte inferiore del corpo presentano \textbf{valvole}, pieghe della
tonaca intima coperte da endotelio, che ricordano e funzionano in maniera simile alle valvole
semilunari.

\rubrica{Fistola artero-venosa}
Comunicazione anomala tra un'arteria e una vena: lesione a carico di un'arteria e di una vena
affiancate. La lesione è solitamente rappresentata da una ferita penetrante, come quella provocata
da un coltello o un proiettile. Può essere congenita (rara) o acquisita.

% ---------------------------------------------------------------------------
\section{Capillari e plessi capillari}

I capillari costituiscono l'unico distretto cardiovascolare ove sono possibili scambi tra sangue
e liquido interstiziale circostante, e quindi cellule e tessuti. Costituiscono una rete a maglie
intricate (plesso capillare, o letto capillare) in cui il sangue scorre a velocità estremamente
bassa, utile per la diffusione o il trasporto attivo transmembrana dei materiali.

Normalmente il sangue procede dai capillari alle vene, per poi far ritorno al cuore. Vi sono 3
eccezioni:
\begin{itemize}
  \item \textbf{ipofisi}: circolazione portale ipotalamo-ipofisaria;
  \item \textbf{fegato}: circolazione portale;
  \item \textbf{rene}: il nefrone presenta un letto capillare arterioso seguito da uno
    arterovenoso.
\end{itemize}

% ---------------------------------------------------------------------------
\section{Circolazione}

\rubrica{La grande circolazione}
Detta anche \textbf{circolazione sistemica}: è la circolazione sanguigna tra il cuore, i tessuti
dell'organismo e di ritorno al cuore. Dal cuore partono le arterie, nelle quali circola il sangue
ricco di ossigeno che raggiunge i tessuti dell'organismo, dove verrà ceduto l'ossigeno ed in
cambio verrà acquistata anidride carbonica. Dai tessuti, il sangue ricco di anidride carbonica
farà ritorno al cuore attraverso le vene.

\rubrica{La piccola circolazione}
Detta anche \textbf{circolazione polmonare}: è la circolazione sanguigna tra il cuore, i polmoni e
di ritorno al cuore. Il sangue ricco di anidride carbonica proveniente dalla circolazione
sistemica viene pompato dal cuore verso i polmoni, dove avviene lo scambio gassoso: viene ceduta
anidride carbonica ed acquistato ossigeno. Il sangue ritorna al cuore, il quale lo pompa nella
circolazione sistemica verso i tessuti dell'organismo.

Il cuore è quindi una pompa formata funzionalmente da 2 metà:
\begin{itemize}
  \item \textbf{destra}, che riceve e spinge sangue deossigenato;
  \item \textbf{sinistra}, che riceve e spinge sangue ossigenato.
\end{itemize}

I ventricoli conferiscono al sangue una consistente pressione: il sistema circolatorio è quindi
organizzato in 2 componenti, una ad alta pressione in partenza dal cuore, e una a bassa pressione
in arrivo al cuore.

% ---------------------------------------------------------------------------
\section{Aorta}

L'aorta si distingue in 3 porzioni: aorta ascendente (dalla quale originano le coronarie), arco
dell'aorta (dal quale originano il tronco brachiocefalico, la carotide comune e la succlavia di
sinistra), aorta discendente (a sua volta suddivisa in 2 porzioni consecutive: toracica ed
addominale). Termina a livello di L4, dividendosi nelle 2 arterie iliache comuni, destra e
sinistra.

\rubrica{Aorta ascendente}
Inizia dalla valvola semilunare aortica, all'altezza della 3\textsuperscript{a} cartilagine
costale di sinistra, margine inferiore. All'origine sono presenti i 3 seni aortici (di Valsalva),
quindi un rigonfiamento detto \textbf{bulbo aortico}. Si porta in alto e a destra, e all'altezza
del margine superiore della 2\textsuperscript{a} cartilagine costale trapassa nell'arco aortico.
Rami: le 2 arterie coronarie, destra e sinistra.

\rubrica{Arco aortico}
Si porta posteriormente, scavalca la biforcazione del tronco polmonare, passa a sinistra della
trachea e poi dell'esofago, passa superiormente al bronco di sinistra, e raggiunge il corpo di
T4. Genera l'arteria anonima (o tronco brachiocefalico), la carotide comune di sinistra e la
succlavia di sinistra. Nel feto, l'arteria polmonare comunica con l'arco aortico mediante il
\textit{condotto di Botallo}.

\rubrica{Aorta discendente}
Porta il sangue direttamente dal cuore al resto del corpo. Si distinguono 2 porzioni: toracica e
addominale; il limite è dato dal passaggio nel diaframma, a livello della dodicesima vertebra
toracica. Da entrambe le porzioni originano rami collaterali che irrorano le pareti (rami
parietali) o i visceri (rami viscerali).

\rubrica{Aorta discendente toracica}
\begin{itemize}
  \item arterie \textbf{bronchiali}: servono bronchi e polmoni;
  \item arterie \textbf{pericardiche}: servono il pericardio;
  \item arterie \textbf{mediastiniche}: servono gli organi mediastinici;
  \item arterie \textbf{esofagee}: servono l'esofago;
  \item arterie \textbf{freniche superiori}: servono il diaframma;
  \item arterie \textbf{intercostali}: servono la parete toracica.
\end{itemize}

\rubrica{Aorta addominale}
Porzione inferiore dell'aorta discendente. Genera:
\begin{itemize}
  \item rami \textbf{posteriori}, parietali, pari;
  \item rami \textbf{laterali}, viscerali, pari (reni, surreni, genitali);
  \item rami \textbf{anteriori}, viscerali, impari.
\end{itemize}

Nel dettaglio:
\begin{itemize}
  \item rami viscerali \textbf{anteriori}:
  \begin{enumerate}
    \item \textbf{tronco celiaco}: 3 rami, al fegato, alla cistifellea, all'esofago, allo stomaco,
      al duodeno, al pancreas, alla milza;
    \item \textbf{mesenterica superiore}: al pancreas e al duodeno, al piccolo intestino e al
      colon;
    \item \textbf{mesenterica inferiore}: al colon terminale e al retto;
  \end{enumerate}
  \item rami viscerali \textbf{laterali}: surrenali (alle ghiandole surrenali), renali (ai reni),
    spermatiche (alle gonadi: ovaie e testicoli);
  \item rami \textbf{parietali}: lombari (al dorso).
\end{itemize}

L'aorta addominale è l'arteria principale della cavità addominale, ultimo tratto dell'aorta e
continuazione dell'aorta toracica: ha una serie di piccoli rami.
\begin{itemize}
  \item \textbf{parietali} (pari): l'arteria frenica inferiore (per il diaframma) e le arterie
    lombari che, in numero di quattro per lato, irrorano i muscoli della parete addominale;
  \item \textbf{viscerali}: l'arteria celiaca, le arterie mesenterica superiore (vascolarizza
    pancreas, duodeno, intestino tenue e crasso) e mesenterica inferiore (porzione terminale del
    colon e retto), l'arteria renale, l'arteria surrenale media e l'arteria genitale.
\end{itemize}

\rubrica{Arterie iliache}
Vasi sanguigni di grandi dimensioni in regione pelvica: a ciascuna arteria iliaca comune fa capo
l'irrorazione dei visceri e delle pareti pelviche, nonché dell'intero arto inferiore
corrispondente.

\begin{itemize}
  \item \textbf{arteria iliaca esterna}: a livello del legamento inguinale termina nell'arteria
    femorale comune, che si divide in un ramo superficiale e uno profondo;
  \item \textbf{arteria iliaca interna}: nasce a livello dell'articolazione sacro-iliaca e si porta
    verso il basso e indietro fino al margine superiore del grande foro ischiatico. Incrocia i
    muscoli grande psoas, piriforme ed il tronco nervoso lombo-sacrale, per poi dividersi in un
    tronco anteriore ed uno posteriore.
\end{itemize}

\rubrica{Aneurisma}
Dilatazione della parete di un'arteria, di una vena o del cuore. Gli aneurismi arteriosi si
manifestano come dilatazioni pulsanti del vaso: la localizzazione più importante è a carico
dell'aorta, nel $75\%$ dei casi colpisce l'aorta addominale. La causa principale è
l'\textbf{aterosclerosi}.

\rubrica{Aterosclerosi}
Malattia infiammatoria cronica delle arterie di grande e medio calibro, che si instaura a causa
dei fattori di rischio cardiovascolare: ipercolesterolemia, diabete, ipertensione, obesità; altre
cause sono di natura infettiva e immunologica. Anatomicamente, la lesione caratteristica è la
\textbf{placca aterosclerotica}, ossia un ispessimento dell'intima (lo strato più interno delle
arterie, rivestito dall'endotelio ed in diretto contatto con il sangue) dovuto principalmente
all'accumulo di materiale lipidico e alla proliferazione del tessuto connettivo.

% ---------------------------------------------------------------------------
\section{Circolazione coronarica}

\begin{itemize}
  \item \textbf{arteria coronaria destra}: segue il solco coronario ed irrora l'atrio destro,
    porzioni di entrambi i ventricoli e parti del sistema di conduzione del cuore (nodo
    senoatriale e nodo atrioventricolare); si dirama nelle arterie marginali e nella
    interventricolare (o discendente posteriore);
  \item \textbf{arteria coronaria sinistra}: irrora atrio e ventricolo sinistro e setto
    interventricolare; si dirama nelle arterie circonflessa ed interventricolare (o discendente
    anteriore).
\end{itemize}

\begin{itemize}
  \item la \textbf{grande vena cardiaca} inizia lungo il solco interventricolare anteriore,
    raccoglie la vena cardiaca posteriore e la vena cardiaca media, e si apre nel seno coronario,
    che sbocca nell'atrio destro vicino allo sbocco della vena cava inferiore;
  \item la \textbf{piccola vena cardiaca} e le vene cardiache anteriori si svuotano nel seno
    coronario e direttamente nell'atrio destro.
\end{itemize}

% ---------------------------------------------------------------------------
\section{Circolazione cerebrale}

\rubrica{Arterie carotidi}
Sono i più grandi tronchi arteriosi del corpo umano: insieme all'arteria vertebrale, la carotide
irrora il SNC e la testa/collo. Ha un primo ramo unico (arteria carotide comune) per poi
suddividersi in arteria carotide interna ed arteria carotide esterna.

Ha diversa origine nel lato destro e nel lato sinistro del corpo: a destra origina dall'arteria
brachiocefalica, mentre a sinistra proviene direttamente dall'aorta. L'arteria carotide comune
risale indicativamente fino all'altezza della laringe, dove si divide in arteria carotide interna
ed arteria carotide esterna.

\rubrica{Poligono di Willis}
Sistema di anastomosi arteriose a pieno canale, presente alla base della scatola cranica: dovrebbe
garantire un'appropriata irrorazione dei tessuti cerebrali. È dato dalla confluenza di 3 arterie
principali:
\begin{itemize}
  \item l'\textbf{arteria basilare}, formata dalla confluenza delle arterie vertebrali destra e
    sinistra (prime collaterali della succlavia);
  \item le \textbf{2 arterie carotidi interne} (destra e sinistra).
\end{itemize}

Il poligono di Willis può essere ricondotto idealmente ad un eptagono avente come lati:
anteriormente le 2 arterie cerebrali anteriori (destra e sinistra), che si uniscono attraverso
l'arteria comunicante anteriore; posteriormente le 2 arterie cerebrali posteriori (destra e
sinistra); tra arteria cerebrale anteriore e posteriore di ogni lato c'è l'arteria comunicante
posteriore.

Sono possibili delle \textbf{varianti}: a livello di origine; di decorso (tortuosità); di calibro;
del sifone carotideo.

Quando la circolazione del sangue nel cervello è bloccata, possono insorgere: ictus, trombosi,
embolo, stenosi.

% ---------------------------------------------------------------------------
\section{Arteria ascellare}

L'\textbf{arteria ascellare}, pari e simmetrica, è la principale arteria del cavo ascellare:
fornisce sangue alla parte superiore del torace, alla spalla ed alla scapola. Costituisce il
prosieguo dell'arteria succlavia.

I suoi rami collaterali sono l'arteria toraco-acromiale, l'arteria sottoscapolare, le toraciche
suprema e laterale, e le circonflesse posteriore ed anteriore dell'omero. Diventa arteria
brachiale al margine inferiore del tendine del muscolo grande pettorale. Ha una lunghezza
complessiva di $11\,cm$ ed un calibro variabile, che oscilla da $0,9\,cm$ alla sua origine a
$0,7\,cm$.

% ---------------------------------------------------------------------------
\section{Il ritorno venoso}

\textbf{Ritorno venoso}: quantità di sangue che ritorna al cuore in un minuto.

La \textbf{curva del ritorno venoso} descrive la relazione tra l'output cardiaco, che dipende da
post-carico, gittata cardiaca, volume sistolico ed \textit{HR}, e le conseguenti variazioni del
flusso del sangue che ritorna al cuore, ovvero del pre-carico. In condizioni fisiologiche il
precarico risulta il principale fattore che limita la massima gittata cardiaca, con un andamento
descritto dalla \textbf{legge di Frank-Starling}. La sovrapposizione della curva della funzione
cardiaca e della curva del ritorno venoso è utilizzata in un modello emodinamico.

\rubrica{Ritorno venoso del cranio}
Le vene che drenano il sangue dagli emisferi cerebrali si gettano in una rete di seni durali,
canali dilatati che percorrono la meninge, di cui i più importanti sono il seno sagittale
superiore, il seno retto (ove si getta la grande vena cerebrale) e il seno cavernoso. I seni
durali danno origine alla vena giugulare interna, che scende lungo il collo parallela all'arteria
carotide interna. Le vene vertebrali scendono nei fori trasversari e defluiscono nelle vene
brachiocefaliche.

\rubrica{Ritorno venoso della testa e del collo}
Le vene superficiali della testa (vene temporali e mascellari) drenano nella vena giugulare
esterna, che scende verso il torace sulla superficie esterna del muscolo sternocleidomastoideo e
sbocca nella vena succlavia; oppure nella vena giugulare interna (v.\ faciale).

\rubrica{Ritorno venoso dagli arti superiori}
\begin{itemize}
  \item \textbf{circolo superficiale}: la vena cefalica risale il lato radiale dell'avambraccio,
    mentre la vena basilica risale il lato ulnare; davanti al gomito passa la vena mediana del
    gomito, da cui solitamente si effettuano i prelievi ematici;
  \item \textbf{circolo profondo}: la vena radiale e la vena ulnare si fondono a formare la vena
    brachiale, che confluisce con la basilica a dare la vena ascellare.
\end{itemize}

\rubrica{Ritorno venoso dagli arti inferiori}
\begin{itemize}
  \item \textbf{circolo superficiale}: la vena grande safena (la più lunga del corpo) risale lungo
    la superficie mediale della gamba e della coscia, e si getta nella vena femorale appena prima
    che penetri nella parete addominale; la vena piccola safena risale lungo la superficie
    laterale della gamba e confluisce nella poplitea;
  \item \textbf{circolo profondo}: le vene tibiale anteriore e posteriore e la vena peroniera si
    riuniscono nella vena poplitea, che prosegue come femorale e va a formare l'iliaca esterna.
\end{itemize}

\rubrica{Ritorno venoso da pelvi e addome}
Le vene iliache esterne si fondono con le iliache interne, che drenano il sangue dagli organi
pelvici, a dare la vena iliaca comune; questa si fonde con la controlaterale davanti a L3 a
formare la \textbf{vena cava inferiore}, che riceve le vene lombari, genitali, renali, surrenali e
freniche.

\rubrica{Vena porta}
Raccoglie il sangue dal tubo gastroenterico, dalla milza e dal pancreas, e lo porta al fegato.
Origina dalle due vene mesenteriche (superiore ed inferiore), dalla vena lienale, dalla vena
cistica e dalle vene gastriche di destra e di sinistra.

% ---------------------------------------------------------------------------
\section{Circolazione venosa}

Si articola in: sistema delle vene cave; sistema delle vene azygos; sistema della vena porta;
circolazione venosa superficiale; circolazione venosa profonda.

La \textbf{vena cava superiore} si forma dalla confluenza dei 2 tronchi brachiocefalici venosi,
destro e sinistro, e riceve sangue anche dalla vena azygos; i tronchi brachiocefalici sono a loro
volta formati dalla confluenza della vena succlavia e della vena giugulare interna.

\rubrica{Vene superficiali del collo}
\begin{itemize}
  \item vena \textbf{giugulare anteriore}: dalla regione sottomandibolare scende verticalmente e
    centralmente per aprirsi nella succlavia;
  \item vena \textbf{giugulare esterna}: dall'angolo della mandibola scende per aprirsi nella
    succlavia.
\end{itemize}

\rubrica{Vena azygos}
Vaso impari asimmetrico, situato nella cavità toracica e addominale: può assumere una grande
importanza in certe situazioni, ad esempio quando la vena cava inferiore risulta ostruita, poiché
il sangue refluo della parte sotto-diaframmatica del corpo può raggiungere ugualmente il cuore
attraverso di essa.

\rubrica{Vene iliache comuni}
Destra (più corta) e sinistra (più lunga): sono i 2 vasi dalla cui convergenza origina la vena
cava inferiore, nella regione pelvica. Originano entrambe dalla confluenza della vena iliaca
esterna con la vena iliaca interna omolaterale, davanti all'articolazione sacro-iliaca
corrispondente.

La \textbf{vena iliaca interna} (o vena ipogastrica) drena il sangue proveniente dalla pelvi
attraverso i suoi rami viscerali e parietali.
